\chapter{Reference: mouse}
\label{ch:refman-mouse}

\begin{aside}
Mouse support is opt-in. Set \code{RunConfig::mouse = true},
add a mouse handler in \code{subscribe}, and clicks /
scrolls / drags arrive as messages.
\end{aside}

\section{Event types}

\begin{itemize}[leftmargin=*]
  \item \code{MouseButton::Left}, \code{Middle},
    \code{Right}.
  \item \code{MouseKind::Press}, \code{Release},
    \code{Drag}, \code{Move}.
  \item Wheel: \code{MouseButton::WheelUp},
    \code{WheelDown}.
\end{itemize}

\section{Working example}

\begin{lstlisting}[language=C++]
#include <maya/maya.hpp>

using namespace maya;
using namespace maya::dsl;

struct MouseDemo {
    struct Model {
        int x = 0, y = 0;
        std::string last = "(none)";
    };
    struct MouseEvt { int x; int y; std::string label; };
    struct Quit {};
    using Msg = std::variant<MouseEvt, Quit>;

    static Model init() { return {}; }

    static auto update(Model m, Msg msg)
        -> std::pair<Model, Cmd<Msg>>
    {
        return std::visit(overload{
            [&](MouseEvt e) {
                m.x = e.x; m.y = e.y; m.last = e.label;
                return std::pair{m, Cmd<Msg>{}};
            },
            [](Quit) {
                return std::pair{Model{}, Cmd<Msg>::quit()};
            },
        }, msg);
    }

    static Element view(const Model& m) {
        return v(
            text(std::format("pos: ({}, {})", m.x, m.y)),
            text(std::format("evt: {}", m.last)),
            blank_,
            t<"click anywhere, q to quit"> | Dim
        ) | pad<1> | border_<Round>;
    }

    static auto subscribe(const Model&) -> Sub<Msg> {
        return Sub<Msg>::batch(
            key_map<Msg>({ {'q', Quit{}} }),
            mouse_map<Msg>([](MouseEvent e) -> Msg {
                std::string lbl;
                switch (e.kind) {
                    case MouseKind::Press:   lbl = "press"; break;
                    case MouseKind::Release: lbl = "release"; break;
                    case MouseKind::Drag:    lbl = "drag"; break;
                    case MouseKind::Move:    lbl = "move"; break;
                }
                return MouseEvt{e.x, e.y, lbl};
            })
        );
    }
};

int main() {
    return run<MouseDemo>({
        .title = "mouse",
        .mouse = true,
    });
}
\end{lstlisting}

\section{Notes}

\begin{itemize}[leftmargin=*]
  \item Coordinates are zero-based; \code{(0, 0)} is the
    top-left cell.
  \item Some terminals require additional config to forward
    mouse events through tmux/screen.
  \item Selection-by-drag in the terminal may conflict with
    \maya{} mouse capture; users hold Shift to override.
\end{itemize}

\section{Recap}

\begin{recap}
\begin{itemize}[leftmargin=*]
  \item Opt in via \code{RunConfig}.
  \item Press, release, drag, move, wheel.
  \item Coordinates in terminal cells.
\end{itemize}
\end{recap}

\section*{Workshop}

\begin{enumerate}[leftmargin=*]
  \item Build a clickable button: highlight on press,
    fire a command on release.
  \item Implement drag-to-scroll on a long list.
  \item Show a tooltip on hover.
\end{enumerate}
